% TOP_PROBLEM: {"id": 7000020, "title": "Geometry of Curves and Surfaces — Problem 5.1", "category_id": 6, "category": {"id": 6, "name": "geometry", "display_name": "Geometry", "description": "Euclidean and non-Euclidean geometry, geometric structures.", "slug": "geometry", "order_index": 6, "created_at": "2026-07-31T15:26:25.671Z"}, "status": "partially_solved"}
% FIRST_POSTED: null
% ENTRY_KIND: statement_only
% Imported from: batch18.tex; UnsolvedMath ID 7000020
\documentclass[11pt]{article}
\usepackage[margin=1in]{geometry}
\usepackage{amsmath,amssymb,mathtools}
\usepackage{fontspec,unicode-math}
\setmainfont{Latin Modern Roman}
\setsansfont{Latin Modern Sans}
\setmathfont{Latin Modern Math}
\usepackage{xurl,hyperref}
\hypersetup{colorlinks=true,urlcolor=blue,linkcolor=blue}
\newcommand{\sref}[2]{\hyperlink{source:#1:#2}{[S#2]}}
\newcommand{\cataloguescope}[1]{\paragraph{Recorded mathematical question.}#1}
\begin{document}
% UnsolvedMath ID 7000020; source catalogue code AMR-069-0020
\setcounter{section}{7000019}
\section{Geometry of Curves and Surfaces — Problem 5.1}\label{7000020}
{\small\sffamily UnsolvedMath ID 7000020 · Geometry}
\subsection{Definitions and mathematical statement}

\paragraph{Shared notation.}$\mathbb N=\{1,2,\ldots\}$, and $\mathbb Z,\mathbb Q,\mathbb R,\mathbb C$ denote the integers, rationals, reals and complex numbers. Any other conventions are specified locally.


A curve is a continuous map $\gamma:[a,b]\to\mathbb R^3$. Its length is
\[L(\gamma)=\sup_{a=t_0<\cdots<t_m=b}\sum_{j=1}^m\lVert\gamma(t_j)-\gamma(t_{j-1})\rVert.\]
It is rectifiable if this is finite, and closed if $\gamma(a)=\gamma(b)$. Its width is
\[w(\gamma)=\min_{u\in S^2}\left(\max_{t\in[a,b]}\langle\gamma(t),u\rangle-\min_{t\in[a,b]}\langle\gamma(t),u\rangle\right),\]
the minimum distance between parallel planes enclosing it. Write $\operatorname{conv}(\gamma)$ for the convex hull of its image. Its inradius is
\[r(\gamma)=\sup\{r>0:\text{some }c+rS^2\subset\operatorname{conv}(\gamma)\text{ is disjoint from }\gamma([a,b])\}.\]
Thus the source imposes disjointness from the curve in addition to containment of the sphere in its convex hull. Consider positive prescribed width or inradius.
\paragraph{Statement recorded in the supplied source.}
Determine the shortest rectifiable curve in $\mathbb R^3$ with a prescribed width, and determine the shortest such curve with a prescribed inradius. In each problem determine also the shortest curve when it is required to be closed. \sref{7000020}{2}
\subsection{Short English statement}
Find the shortest space curve of a given width or inradius, and solve both problems also for closed curves.
\subsection{Sources}
\begin{description}
\item[\hypertarget{source:7000020:1}{S1}] ULAM.AI and the UnsolvedMath Contributors, \emph{UnsolvedMath: A Curated Collection of Open Mathematics Problems}, record 7000020 (AMR-069-0020). CC BY 4.0. \url{https://huggingface.co/datasets/ulamai/UnsolvedMath}.
\item[\hypertarget{source:7000020:2}{S2}] Mohammad Ghomi, \emph{The length, width, and inradius of space curves}, arXiv:1605.01144v3 (2018), Introduction, pp. 2--4; Section 2, definition of curve and length, p. 4. \url{https://arxiv.org/pdf/1605.01144v3}.
\end{description}
\label{end:7000020}
\end{document}
